\section{Literature Review and Multi-dimensional Taxonomy}
\label{sec:lit-review}

Exact subgraph matching has been studied for more than one research generation, and representative progress has appeared from both the database and graph-algorithm communities. Prior work does not advance along a single axis: some methods emphasize aggressive candidate-space reduction before backtracking begins, some optimize the search order to expose contradictions earlier, some reuse dead-end information to avoid repeating equivalent failures, and others redesign the whole execution pipeline around candidate representation and set operations~\cite{he2008graphql,han2019daf,sun2020study,sun2021rapidmatch,kim2022veq,arai2023gup}. Therefore, a meaningful literature review for this course project should not merely list papers chronologically. Instead, it should provide a \emph{multi-dimensional taxonomy} that explains how different methods trade off filtering power, search overhead, and implementation complexity, and how the reproduced GuP method fits into this larger design space.

\subsection{Overview of Prior Work}

Prior work on exact subgraph matching is highly diverse, so in this report we summarize it with a compact multi-dimensional taxonomy instead of a purely chronological list. The goal is to highlight the dominant optimization ideas behind representative methods and to clarify how GuP relates to them.

\subsection{Taxonomy Dimensions}

We classify representative work using four main dimensions.

\begin{itemize}
  \item \textbf{Candidate filtering}: how aggressively a method shrinks the domain of each query vertex before or during the search. Typical signals include labels, degrees, neighborhood signatures, and local structural consistency. Stronger filtering reduces the branching factor, but it may require heavier preprocessing or more complex auxiliary data structures~\cite{he2008graphql,sun2020study}.

  \item \textbf{Matching order}: how the method chooses the next query vertex to expand. Since exact matching is usually realized as state-space search, the order of expansion strongly affects whether contradictions appear early or late. Static heuristics are simple and cheap, while adaptive order selection can better react to the current partial embedding but also introduces additional control overhead~\cite{han2019daf,sun2020study}.

  \item \textbf{Conflict reuse / failing sets}: whether the method memorizes information learned from dead ends and reuses it to skip equivalent or dominated branches later. This dimension is especially important when many branches fail for similar reasons; in such cases, remembering the ``shape'' of failure may be more valuable than further tightening local consistency checks~\cite{han2019daf,arai2023gup}.

  \item \textbf{System-level acceleration}: whether the method relies on heavier indexing, candidate-space engineering, specialized set operations, or a holistic execution framework beyond local pruning logic. This dimension separates algorithms that are conceptually lightweight but search-aware from methods whose performance depends heavily on broader system design~\cite{he2008graphql,sun2021rapidmatch}.
\end{itemize}

These dimensions are complementary rather than mutually exclusive. A strong exact matcher often combines several of them, but different papers prioritize different components. Accordingly, our taxonomy is intended to show \emph{dominant emphasis} rather than force each method into a single exclusive box.

\begin{table*}[t]
  \caption{A multi-dimensional taxonomy of representative exact subgraph matching studies.}
  \label{tab:taxonomy}
  \centering
  \scriptsize
  \setlength{\tabcolsep}{2.5pt}
  \renewcommand{\arraystretch}{1.05}
  \begin{tabular}{p{0.11\textwidth}ccccp{0.18\textwidth}p{0.13\textwidth}}
    \toprule
    Study / Method & Cand. filt. & Order & Conflict reuse & System support & Main emphasis & Relation to GuP \\
    \midrule
    GraphQL~\cite{he2008graphql} & Strong & Moderate & No & Yes & Signature-based pruning and graph access methods for reducing candidate space early & Earlier candidate-space-oriented line \\
    DAF~\cite{han2019daf} & Strong & Strong & Strong & Moderate & Dynamic programming, adaptive matching order, and failing sets in one exact matcher & Closest classical conflict-aware reference \\
    In-depth study~\cite{sun2020study} & Covers many & Covers many & Covers many & Covers many & Comparative analysis showing which components dominate performance across workloads & Evaluation vocabulary and taxonomy support \\
    RapidMatch~\cite{sun2021rapidmatch} & Strong & Strong & Limited & Strong & Holistic high-performance query processing instead of a single pruning trick & Stronger system emphasis than GuP \\
    VEQ~\cite{kim2022veq} & Moderate & Moderate & Limited & Moderate & Equivalence-aware reduction of redundant matching work & Alternative redundancy-reduction line \\
    GuP~\cite{arai2023gup} & Moderate & Moderate & Strong & Moderate & Guard-based pruning through reservation and nogood mechanisms & Our reproduction target \\
    \bottomrule
  \end{tabular}
\end{table*}

\subsection{Representative Research Lines}

\paragraph{Candidate-space-oriented filtering.}
An influential early database line focuses on reducing the candidate space as much as possible before expensive enumeration begins. GraphQL is a representative example: it integrates signature-based pruning and graph access methods so that many impossible matches can be removed early, before full backtracking search is triggered~\cite{he2008graphql}. This line established an important intuition for later work: exact matching performance is often determined by how small and how accurate the candidate domains are. In practice, however, such methods can still struggle when local signatures are insufficiently discriminative, because substantial ambiguity remains and must be resolved by downstream search.

\paragraph{Integrated search optimization.}
A second line treats exact subgraph matching explicitly as a search problem whose major components should be optimized together. DAF is representative of this direction because it combines dynamic programming, adaptive matching order, and failing sets within one exact matcher~\cite{han2019daf}. The importance of DAF in our review is twofold. First, it shows that search order is not a secondary heuristic but a first-class performance factor: different orderings can expose conflicts at very different depths. Second, it demonstrates that failure information can be summarized and reused, rather than discarded after each dead end. Compared with earlier candidate-space-oriented methods, this line moves more of the optimization effort from static preprocessing to dynamic search control.

\paragraph{Comparative understanding of in-memory systems.}
The field has also developed a more reflective line of work that compares existing exact matchers systematically instead of proposing only one new algorithm. The in-depth study by Sun and Luo is especially valuable in this regard~\cite{sun2020study}. Rather than attributing performance to a single pruning idea, it analyzes how candidate generation, matching order, intermediate representation, and implementation decisions interact across workloads. For our course project, this study is useful because it provides a vocabulary for discussing why a method is fast or slow. It also supports our decision to report not only runtime, but also search-space indicators such as the number of partial mappings.

\paragraph{Holistic execution and system support.}
Another research line pushes exact matching toward a more holistic execution-engine perspective. RapidMatch is representative here: instead of emphasizing only one local pruning rule, it redesigns subgraph query processing around the coordinated treatment of candidate generation, query execution, and intermediate-state handling~\cite{sun2021rapidmatch}. This line is important because it reminds us that state-of-the-art performance may rely not only on better pruning logic, but also on better organization of the entire execution pipeline. From a reproduction perspective, however, such methods are often harder to emulate faithfully in a course project, since low-level engineering choices contribute substantially to the final speedup.

\paragraph{Redundancy and equivalence reduction.}
A closely related but conceptually distinct line seeks to reduce redundant search by recognizing when different branches correspond to equivalent or near-equivalent computational states. VEQ is a representative method in this category~\cite{kim2022veq}. Its contribution is useful for our taxonomy because it shows that repeated work can be reduced not only by stronger filtering or better failure memorization, but also by detecting equivalence relationships among search states. Compared with DAF-style failing sets, this direction is less about recording explicit dead ends and more about compressing duplicated work more broadly.

\paragraph{Relation among these lines.}
Taken together, these studies suggest that exact subgraph matching has evolved along several complementary directions rather than a single dominant path. Candidate-space methods try to shrink the search problem early; search-centric methods optimize how the search tree is traversed; conflict-reuse methods try to avoid rediscovering the same failures; and system-oriented methods improve the execution pipeline as a whole. This perspective is more useful for our report than a simple chronological list, because it directly supports the later analysis of GuP's reservation guard and nogood guard.

\subsection{Positioning of GuP in the Taxonomy}

Under the above taxonomy, GuP is best understood as a \emph{search-based exact matcher with dynamic state-aware pruning}. Its core contribution is not the heaviest pre-search filtering, nor the most system-heavy execution engine, but the introduction of guards that evaluate whether a candidate extension is likely to make the remaining search infeasible~\cite{arai2023gup}. In particular, the reservation guard propagates future injectivity pressure earlier in the search, while the nogood guard attempts to reuse information from previously discovered dead ends. Therefore, GuP lies between classical candidate-space methods and stronger failing-set-style conflict-reuse methods: it still operates inside a backtracking enumeration framework, but it pushes pruning decisions closer to the dynamic search state.

This position makes GuP especially suitable for a course-project reproduction. Compared with GraphQL, GuP is less dependent on heavyweight offline candidate-space summaries; compared with DAF, it exposes its pruning ideas through relatively modular guards; compared with RapidMatch, it is less tied to full-stack system engineering; and compared with VEQ, its pruning contribution is easier to align with the course-required statistics on partial mappings and filtered partial mappings~\cite{he2008graphql,han2019daf,sun2021rapidmatch,kim2022veq,arai2023gup}. In other words, GuP offers a favorable balance between research novelty, conceptual clarity, ablation friendliness, and implementability.

\subsection{Implications for the Rest of This Report}

The literature review leads to three concrete decisions for the rest of our report. First, we treat GuP as part of the broader conflict-aware exact matching line rather than as an isolated technique. Second, we evaluate our prototype not only by end-to-end runtime but also by search-space indicators that reflect how the guards affect enumeration. Third, when discussing limitations, we distinguish algorithmic limitations from system-level ones: if our Python reproduction underperforms highly optimized prior systems, this does not automatically invalidate the pruning ideas themselves. This distinction is important for interpreting experimental results fairly in a course-project setting.
